Let $\sigma$ denote a random permutation from
the symmetric group $\cS_n$, each permutation being equally
likely\footnote{Random permutations sampled
from certain other distributions have been studied, e.g. \cite{ABTbook}, but we will not discuss these here.}.
We denote by $\PR_n$ and $\E_n$ the probability and expectation
with respect to a uniform random $\sigma\in \cS_n$.
Often, the subscript $n$ will be omitted if it is clear from the context.
We denote the type (or cycle type) of $\sigma$ by
\[
(C_1(\sigma), C_2(\sigma),\ldots, C_n(\sigma)),
\]
where
$C_j(\sigma)$ is the number of cycles of length $j$ in $\sigma$.
More generally, for any subset $I$ of $[n]=\{1,\ldots,n\}$, we let $C_I(\sigma)$
be the number of cycles whose lengths lie in the set $I$.
For brevity, we write $C(\sigma)$ for the total number of
cycles in $\sigma$.
